\addcontentsline{toc}{subsubsection}{Solution 53. When Even Numbers Become Smaller}

\begin{dlsummary}[title={\bfseries Solution 53. When Even Numbers Become Smaller},breakable]

\begin{proof}
The unique solution is
\[
\boxed{
f(0)=0,\qquad
f\!\left(2^k\frac ab\right)=2^{-k},
}
\]
where \(k\in\mathbb Z\), and \(a,b\) are odd integers with \(b>0\).

In other words, only the exponent of \(2\) in a nonzero rational number matters.

\paragraph{1. Establish the basic identities.}

Since
\[
f(2)=f(2)f(1)>0,
\]
we have \(f(1)=1\). Consequently, for \(x\ne0\),
\[
f(x)f(1/x)=1,
\]
so \(f(x)>0\) and
\[
f(1/x)=\frac1{f(x)}.
\]
Also, \(f(-1)^2=f(1)=1\), so nonnegativity gives \(f(-1)=1\). Thus
\[
f(-x)=f(x).
\]

\paragraph{2. Use binary expansions to bound all integers.}

Write a positive integer \(n\) in binary:
\[
n=\varepsilon_0+\varepsilon_1 2+\cdots+\varepsilon_s2^s,
\qquad \varepsilon_j\in\{0,1\}.
\]
Repeated subadditivity gives
\[
f(n)\le\sum_{j=0}^s\varepsilon_j f(2^j)
=\sum_{j=0}^s\varepsilon_j2^{-j}<2.
\]
Therefore
\[
\boxed{f(n)<2\qquad(n\in\mathbb Z_{>0}).} \tag{1}
\]

\paragraph{3. Sharpen the bound by taking powers.}

Apply (1) to \(n^r\). Multiplicativity yields
\[
f(n)^r=f(n^r)<2
\]
for every positive integer \(r\).

If \(f(n)>1\), its powers would eventually exceed \(2\). Hence
\[
\boxed{f(n)\le1\qquad(n\in\mathbb Z_{>0}).} \tag{2}
\]

\paragraph{4. Determine the values at odd integers.}

Let \(n\) be a positive odd integer. For every positive integer \(r\),
\[
j=\frac{n^r-1}{2}
\]
is a nonnegative integer. Since \(1=n^r-2j\), subadditivity, evenness, and (2) imply
\[
1=f(1)\le f(n^r)+f(2j)
=f(n)^r+\frac12f(j)
\le f(n)^r+\frac12.
\]
Thus
\[
f(n)^r\ge\frac12
\]
for every \(r\). If \(f(n)<1\), its powers would tend to zero. Together with (2), this proves
\[
\boxed{f(n)=1\qquad(n\text{ odd}).}
\]

\paragraph{5. Determine all rational values and verify.}

Every nonzero rational has a unique representation
\[
x=2^k\frac ab
\]
with \(a,b\) odd and \(b>0\). Multiplicativity now forces
\[
f(x)=f(2)^k\frac{f(a)}{f(b)}=2^{-k}.
\]

Conversely, define \(f\) by this formula. Multiplication adds the exponents of \(2\), so multiplicativity holds.

For nonzero \(x,y\), write
\[
x=2^k u,\qquad y=2^\ell v,\qquad k\le\ell,
\]
where \(u,v\) have odd numerators and denominators. Then
\[
x+y=2^k\bigl(u+2^{\ell-k}v\bigr).
\]
Unless the sum is zero, the expression in parentheses has a nonnegative exponent of \(2\). Therefore
\[
f(x+y)\le2^{-k}=\max\{f(x),f(y)\}\le f(x)+f(y).
\]
Zero cases are immediate, completing the verification.
\end{proof}

\paragraph{Takeaways.}
\begin{itemize}
\item **A small value at \(2\) makes binary expansions powerful:** They produce a uniform bound independent of the integer.
\item **Powers sharpen uniform bounds:** From \(f(n^r)<2\), multiplicativity forces \(f(n)\le1\).
\item **Parity supplies the reverse bound:** The identity \(n^r-2j=1\) prevents odd integers from having values below \(1\).
\item **Stronger inequalities can emerge:** The solution satisfies the maximum bound \(f(x+y)\le\max\{f(x),f(y)\}\), although only a sum bound was assumed.
\end{itemize}

\end{dlsummary}
